\section{Uniform Dulac estimates}\label{sec:third-derivatives}
The return comparisons require Dulac derivative estimates that remain
uniform as the characteristic ratio passes through one. We construct
common section coordinates on compact hyperbolic saddle families and
bound the logarithmic curvature and its Euler derivative.
All source derivatives keep the field parameters fixed.

\begin{lemma}\label{lem:dulac-jets}
Fix an integer $J\ge4$. Let $p$ range over a compact subset of a
finite-dimensional parameter manifold covered by finitely many charts.
Assume that the vector field, its saddle branch, and the embeddings of two
oriented transverse sections are jointly $C^{J+12}$ in phase, section, and
chart parameters, with common bounds. Suppose that the saddle eigenvalues
are $\lambda_u(p)>0$ and $-\lambda_s(p)<0$, with
$\lambda_u,\lambda_s\ge\lambda_*>0$, and that
$r(p)=\lambda_s(p)/\lambda_u(p)$ lies in
$[r_-,r_+]\Subset(0,\infty)$. After shortening the sections uniformly, the
local coordinates and section changes may be chosen jointly $C^J$, and the
Dulac map $D_p$ satisfies, for $0<s\le s_0$,
\begin{equation}\label{eq:dulac-third}
 D_p'>0,\qquad D_p'\le Ks^{r(p)-1},\qquad
 \frac{sD_p''}{D_p'}=r(p)-1+E_D(s,p),\qquad
 |E_D|+|sE_D'|\le Ks^\nu
\end{equation}
with common $K,s_0,\nu>0$. Primes and $E_D'$ denote source derivatives at
fixed $p$. The estimates include $r=1$; when $r$ varies, no joint $C^J$
extension of $D_p$ to $s=0$ is asserted.
\end{lemma}
\begin{proof}
Put $N=J+12$. The hypotheses supply common $C^N$ phase--parameter bounds
after one common phase shrinking. Cover the compact
parameter set by finitely many charts on which the two eigenline bundles
have $C^{N-1}$ frames with common bounds. Translate the saddle to the
origin and use these frames on each chart to obtain
\begin{equation}\label{eq:dulac-hyperbolic-chart}
 \dot\xi=\lambda_u(p)\xi+F_u(\xi,\eta,p),\qquad
 \dot\eta=-\lambda_s(p)\eta+F_s(\xi,\eta,p),
\end{equation}
where $F(0,p)=D_{(\xi,\eta)}F(0,p)=0$.  Compact hyperbolicity and the
finite cover give a common $\lambda_*>0$ such that
$\lambda_u,\lambda_s\ge\lambda_*$, common $C^{N-1}$ bounds for the
coefficients, and a common bound for the frames and their inverses.  The
loss of one derivative here accounts for differentiating the field when
forming its spectral projections.

To obtain common separatrix jets in phase and parameters, use the
Lyapunov--Perron equation for stable trajectories of
\eqref{eq:dulac-hyperbolic-chart}. For an initial stable coordinate $a$,
it is
\begin{equation}\label{eq:dulac-lyapunov-perron}
 \begin{aligned}
  \xi(t)&=-\int_t^\infty e^{\lambda_u(p)(t-q)}
                  F_u(\xi(q),\eta(q),p)\,dq,\\
  \eta(t)&=e^{-\lambda_s(p)t}a
       +\int_0^t e^{-\lambda_s(p)(t-q)}
                  F_s(\xi(q),\eta(q),p)\,dq .
 \end{aligned}
\end{equation}
In a fixed exponential trajectory norm with weight below $\lambda_*$,
a common cutoff makes this operator a contraction with constant
$q_0<1$, uniformly on each parameter chart. Differentiating its
fixed-point equation gives
\[
 (I-D_w\mathcal T_p)D^m w=\mathcal R_m,\qquad
 \|(I-D_w\mathcal T_p)^{-1}\|\le(1-q_0)^{-1},
\]
where $\mathcal R_m$ uses only lower trajectory jets and coefficients
through order $m$. Kernel parameter derivatives add polynomial factors
in time, absorbed by the spectral gap. Induction through $N-1$ gives
common joint stable-graph jets; time reversal gives the unstable graph.
Take finite minima of the chart radii. Local frames suffice, without
a global orientation of either eigenline bundle.

Straighten the two invariant graphs successively.  Invariance and
Hadamard division then give
\[
 \dot x=xA(x,y,p),\qquad \dot y=-yC(x,y,p),
 \qquad A(0,0,p)=\lambda_u(p),\quad C(0,0,p)=\lambda_s(p).
\]
After one common shrinking, $A>0$; division by this positive time factor
puts the equations in the orbital form
\[
 \dot x=x,\qquad \dot y=-yR(x,y,p),\qquad
 R(0,0,p)=r(p)=\lambda_s(p)/\lambda_u(p).
\]
Absorb the pure-axis terms by one-dimensional coordinate changes.
First put $y_1=q(x)y$, where
\[
 q(x)=\exp\int_0^x\frac{R(t,0)-r}{t}\,dt.
\]
Then $\dot y_1=-y_1R_1(x,y_1)$ with $R_1(x,0)=r$.
The further change
\[
 y_2=\phi(y_1),\qquad
 \phi(v)=v\exp\int_0^v\frac{r/R_1(0,t)-1}{t}\,dt
\]
makes the coefficient equal to $r$ on both axes. Each quotient has
a removable singularity and common finite derivatives.  The remaining
coefficient difference vanishes on both axes, so two successive Hadamard
factorizations make it divisible by $xy$.  Dropping subscripts, the
resulting local orbital form is
\[
 \dot x=x,\qquad \dot y=-y[r+xyB(x,y,p)],
\]
with common finite derivatives of $B$; the mixed term is retained,
including at resonance. The derivative orders at each step are
\[
\begin{array}{@{}lcc@{}}
\text{operation} & \text{input order} & \text{output order}\\
\hline
\text{spectral projections and eigenframes} & J+12 & J+11\\
\text{two graph straightenings and push-forwards} & J+11 & J+9\\
\text{two initial Hadamard divisions by the axes} & J+9 & J+7\\
\text{two one-variable quotient changes and push-forwards} & J+7 & J+3\\
\text{two final Hadamard divisions producing }xyB & J+3 & J+1\\
\text{regular section hit and coordinate transfer} & J+1 & J
\end{array}
\]
Every division is applied only after the required vanishing identity has
been established. The coordinate Jacobians and the transverse event
derivatives are bounded away from zero on the finite cover, so inversion
costs no additional derivative. Thus $N=J+12$ leaves the coordinates,
their inverses, $B$, and the section maps jointly $C^J$, with common bounds.

We now estimate the Dulac map in these coordinates; see also
\cite{MarinVilladelprat2020}. Fix $p$ and suppress it from the notation.
For incoming $(s,y_0)$ and outgoing $x=x_0$, factor out the saddle power by
writing $y=(s/x)^rw$, $w(s;s)=y_0$, and
$\mathcal E=s\partial_s$. Then
\[
 w_x=-(s/x)^rw^2B(x,(s/x)^rw,p).
\]
Choose the section sizes so that $y_0/2\le w\le2y_0$.
For $0<\nu<\min(1,r_-)$, reduced once if necessary,
\[
 \int_s^{x_0}(s/x)^r(1+|\log(x/s)|^j)\,dx\le K_js^\nu
 \quad(0\le j\le3).
\]
The first three Euler differentiations give triangular variational
equations whose coefficients and inhomogeneous terms are bounded by
multiples of $(s/x)^r$. The initial point also moves with $s$, so its
differentiated values satisfy
\[
 (\mathcal E^jw)(s;s)
 =s\frac d{ds}(\mathcal E^{j-1}w)(s;s)
                  -s(\partial_x\mathcal E^{j-1}w)(s;s).
\]
They are $O(s)$, successively for $j=1,2,3$. Gronwall's inequality
therefore gives $|w-y_0|+\sum_{j=1}^3|\mathcal E^jw|\le Ks^\nu$.
At $x=x_0$ let $W=w(x_0;s)$ and $M=rW+\mathcal EW\asymp1$.
Since $D'=x_0^{-r}s^{r-1}M$,
\[
 E_D=\frac{\mathcal EM}{M},\qquad
 \mathcal E^2M=r\mathcal E^2W+\mathcal E^3W,
\]
which proves~\eqref{eq:dulac-third} in each chart.  On chart overlaps,
the source and target coordinates differ by positive regular germs with
common $C^J$ bounds.  The chain rule changes the displayed logarithmic
jet, and its Euler derivative, by $O(s)+O(D(s))$.  Decreasing
$\nu<\min(1,r_-)$ once, these terms are absorbed in $Ks^\nu$.  Taking
the maximum of the finitely many constants proves the uniform assertion.
No exact analytic linearization or nonresonance hypothesis was used.
\end{proof}
